
\section{8. Analysis}

\subsection{8.1 Where Does Primal Attention Win?}

Primal Attention provides three benefits:

\begin{table}[H]
\centering
\begin{tabular}{ll}
\toprule
\textbf{Benefit} & \textbf{Mechanism} \\
\midrule
**Attention memory reduction** & $L \times L \to L \times m$ attention matrix \\
**Inference speedup at long $L$** & $\mathcal{O}(L m)$ instead of $\mathcal{O}(L^2)$ \\
**VRAM reduction** & Fewer intermediate activations \\
\bottomrule
\end{tabular}
\end{table}

At $L = 128$, the overhead of computing $F_1, F_2, F_3$ dominates; at 
$L = 2048$, the $\mathcal{O}(L^2)$ term dominates baseline and Primal 
wins.

\subsection{8.2 Where Does It Lose?}

\subsubsection{8.2.1 Landmark Selection Bias}

Segment pooling is \textbf{content-agnostic}: it treats all positions equally. 
If important content is concentrated in a few positions, pooling dilutes it.

\textbf{Figure 8.1} shows attention weight distributions. For layers where 
attention concentrates on a few tokens, Primal Attention's approximation 
is worse.

\subsubsection{8.2.2 Asymmetry Overhead}

The asymmetric formulation computes three factors ($F_1, F_2, F_3$) instead 
of one. This is $3 \times$ the matmul count, offsetting the memory savings 
at short $L$.

\subsubsection{8.2.3 Softmax Over Approximate Scores}

Primal Attention applies softmax to $\tilde{A}$, which is not a proper 
attention matrix. The softmax of an approximate matrix may amplify errors 
compared to softmax of the exact matrix.

\subsection{8.3 Ablation: Landmark Count}

\textbf{Figure 8.2} shows accuracy vs $m$. The curve has three regimes:

\begin{itemize}
\item \textbf{$m < 32$:} sharp degradation
\end{itemize}
\begin{itemize}
\item \textbf{$32 \le m \le 128$:} plateau
\end{itemize}
\begin{itemize}
\item \textbf{$m > 128$:} diminishing returns
\end{itemize}

\subsection{8.4 Ablation: Pseudoinverse Method}

We compare three methods for computing $F_2^{-1}$:

\begin{table}[H]
\centering
\begin{tabular}{lll}
\toprule
\textbf{Method} & \textbf{Accuracy} & \textbf{Time} \\
\midrule
SVD pinv (exact) & {{SVD\_ACC}} & 1.0$\times$ \\
Schulz from zero & {{SCHULZ0\_ACC}} & 0.8$\times$ \\
Schulz from SVD pinv & **{{D\_ACC}}** & 1.2$\times$ \\
\bottomrule
\end{tabular}
\end{table}

\textbf{Observation.} Initializing from SVD pinv is essential for stability.

\subsection{8.5 Causality Analysis}

For decoder models, we apply causal masking to $\tilde{A}$. However, 
$F_2$ is computed from all landmarks, which see the full sequence. This 
means the approximation is not strictly causal.

\textbf{Effect on accuracy.} On GPT-2, causal Primal Attention drops 
\textbf{{{CAUSAL\_DROP}}\%} compared to bidirectional Primal Attention on the 
same weights.

\textbf{Impact.} For classification tasks (SST-2), this is tolerable. For 
generation tasks, it may be more problematic.

\subsection{8.6 Numerical Stability}

\subsubsection{8.6.1 Pseudoinverse Conditioning}

The condition number of $F_2$ can be large when landmarks are 
insufficient. We regularize via:

\[
F_2^{-1} = \text{pinv}(F_2, \text{rcond}=10^{-4})
\tag{8.1}
\]

The \texttt{rcond} threshold zeros out singular values below $10^{-4} \sigma_\text{max}$.

\subsubsection{8.6.2 FP16 vs FP32}

All attention factors are computed in FP32 for stability. The final softmax 
is cast back to the model's dtype.

\subsubsection{8.6.3 Schulz Stability}

The Schulz iteration diverges if $\|I - F_2 X_0\| > 1$. Starting from 
$X_0 = F_2^+$ gives $\|I - F_2 X_0\| \approx 0$, ensuring convergence.

\subsection{8.7 Per-Layer Analysis}

\textbf{Table 8.1: Reconstruction error per layer (DistilBERT).}

\begin{table}[H]
\centering
\begin{tabular}{lll}
\toprule
\textbf{Layer} & \textbf{$F_2$ condition} & \textbf{Approx error} \\
\midrule
0 & {{D\_L0\_COND}} & {{D\_L0\_ERR}} \\
1 & {{D\_L1\_COND}} & {{D\_L1\_ERR}} \\
2 & {{D\_L2\_COND}} & {{D\_L2\_ERR}} \\
3 & {{D\_L3\_COND}} & {{D\_L3\_ERR}} \\
4 & {{D\_L4\_COND}} & {{D\_L4\_ERR}} \\
5 & {{D\_L5\_COND}} & {{D\_L5\_ERR}} \\
\bottomrule
\end{tabular}
\end{table}

\textbf{Observation.} Later layers have larger condition numbers and worse 
approximation.

\subsection{8.8 Comparison to Nystr mformer}

Nystr mformer uses symmetric landmarks (same for $Q, K$). Primal Attention 
allows asymmetric. Comparison:

\begin{table}[H]
\centering
\begin{tabular}{lll}
\toprule
\textbf{Method} & \textbf{Accuracy} & \textbf{Note} \\
\midrule
Nystr mformer (sym) & {{NYST\_ACC}} & Baseline \\
**Primal (asym)** & **{{D\_ACC}}** & +{{ASYM\_GAIN}}\% \\
\bottomrule
\end{tabular}
\end{table}

\textbf{Observation.} Asymmetry gives a small but consistent gain.

\subsection{8.9 Summary of Insights}

1. \textbf{Asymmetry helps:} +{{ASYM\_GAIN}}\% over symmetric Nystr m.
2. \textbf{$m = 64$ is optimal} for $L = 128$.
3. \textbf{Schulz from SVD pinv} is best.
4. \textbf{Causality is approximate.}
5. \textbf{Long $L$ benefits most.}